\documentclass[11pt,a4paper]{article}
\usepackage[margin=1in]{geometry}
\usepackage[T1]{fontenc}
\usepackage[utf8]{inputenc}
\usepackage{lmodern}
\usepackage{microtype}
\usepackage{amsmath,amssymb,amsthm,mathtools,bm}
\usepackage{booktabs,array,longtable}
\usepackage[hidelinks]{hyperref}
\usepackage{enumitem}
\usepackage{physics}
\usepackage{setspace}
\setstretch{1.06}

% SST macro prelude
\newcommand{\swirlarrow}{\mathchoice{\mkern-2mu\scriptstyle\boldsymbol{\circlearrowleft}}{\mkern-2mu\scriptstyle\boldsymbol{\circlearrowleft}}{\mkern-2mu\scriptscriptstyle\boldsymbol{\circlearrowleft}}{\mkern-2mu\scriptscriptstyle\boldsymbol{\circlearrowleft}}}
\newcommand{\swirlarrowcw}{\mathchoice{\mkern-2mu\scriptstyle\boldsymbol{\circlearrowright}}{\mkern-2mu\scriptstyle\boldsymbol{\circlearrowright}}{\mkern-2mu\scriptscriptstyle\boldsymbol{\circlearrowright}}{\mkern-2mu\scriptscriptstyle\boldsymbol{\circlearrowright}}}
\newcommand{\vswirl}{\mathbf{v}_{\swirlarrow}}
\newcommand{\vswirlcw}{\mathbf{v}_{\swirlarrowcw}}
\newcommand{\SwirlClock}{S_t^{\swirlarrow}}
\newcommand{\SwirlClockcw}{S_t^{\swirlarrowcw}}
\newcommand{\rhoF}{\rho_{\!f}}
\newcommand{\rhoE}{\rho_{\!E}}
\newcommand{\rhoM}{\rho_{\!m}}
\newcommand{\rhocore}{\rho_{\text{core}}}
\newcommand{\rc}{r_c}

\title{Toward a Canonical Atomic Torsion-Bridge in Swirl-String Theory\\
\large A Minimal Field--Source--Atom Model with a Provisional Frozen v0.1 Coupling}
\author{Omar Iskandarani}
\date{March 2026}

\begin{document}
\maketitle

\begin{abstract}
This note extracts the minimal reusable core from three prior Swirl-String Theory (SST) drafts and assembles it into a single atomic bridge model aimed at possible canonical use. The retained ingredients are: (i) a propagating transverse torsion/shear field represented by a divergence-free vector potential, (ii) a rotating-fluid source law that converts axial displacement gradients into vertical vorticity and azimuthal swirl, and (iii) a branch-resolved atomic Hamiltonian in which the two Swirl-Clock sectors $\SwirlClock$ and $\SwirlClockcw$ provide the internal two-valued branch. The resulting model is conservative by construction: the ordinary Coulombic atomic limit is retained, electron screening is handled at the effective mean-field level, and the new SST content is confined to a clock-sector coupling driven by a torsion amplitude field. A hydrogenic limit, a many-electron closure rule, a provisional parameter freeze that removes the free branch couplings, and a small set of falsifiable branch-sensitive predictions are stated explicitly. The intentionally excluded pieces are equally important: the ultra-fast compressional branch, skyrmion-emission analogies, and magneto-optical extensions are left outside the canonical core until independently secured.
\end{abstract}

\section{Aim and canonical constraints}
The objective is not to replace the successful mathematics of atomic physics, but to reinterpret it through a minimal SST field theory while preserving its tested limits. The required observable layer is therefore retained as
\begin{equation}
\label{eq:state-labels}
\text{electron state}=(n,\ell,m_\ell,\sigma),
\qquad
\sigma\in\{\swirlarrow,\swirlarrowcw\},
\end{equation}
with the intended low-energy identification that the SST branch label $\sigma$ reproduces the observed two-valued internal sector ordinarily encoded by $m_s=\pm \tfrac12$.

A minimally canonical SST atomic equation must satisfy four constraints:
\begin{enumerate}[label=(C\arabic*)]
    \item recover the hydrogenic spectrum in the one-electron weak-coupling limit;
    \item preserve the shell-state counting $N_{n\ell}=2(2\ell+1)$ and $N_n=2n^2$;
    \item admit a closed many-electron mean-field interpretation via a self-consistent density $\rho$;
    \item isolate new SST content to explicit, testable couplings rather than hidden redefinitions.
\end{enumerate}

\section{Retained ingredients from the earlier drafts}
Three ingredients are retained.

\paragraph{(i) Propagating torsion projection.} The earlier photon/torsion draft identifies an electromagnetic-strength two-form as a projection of torsion,
\begin{equation}
\label{eq:Ftorsion}
F \equiv \kappa\,u_a T^a,
\end{equation}
in a teleparallel setting, and models the propagating radiation sector as a transverse director/shear wave with
\begin{equation}
\label{eq:uA}
\mathbf{u}=\partial_t \mathbf{A},
\qquad
\nabla\cdot \mathbf{A}=0,
\end{equation}
obeying the massless wave equation.

\paragraph{(ii) Rotating-fluid source law.} The rotating-cylinder drafts provide a concrete source chain in a uniformly rotating incompressible fluid. For axisymmetric perturbations,
\begin{equation}
\label{eq:wz}
\partial_t \omega_z = 2\Omega\,\partial_z w,
\qquad
w=\partial_t \xi,
\end{equation}
which integrates from rest to
\begin{equation}
\label{eq:wzxi}
\omega_z(r,z,t)=2\Omega\,\partial_z \xi(r,z,t).
\end{equation}
The azimuthal swirl then follows kinematically from
\begin{equation}
\label{eq:utheta}
\omega_z = \frac{1}{r}\partial_r(r u_\theta)
\quad\Longrightarrow\quad
u_\theta(r,z,t)=\frac{1}{r}\int_0^r \omega_z(r',z,t)\,r'\,dr'.
\end{equation}
This yields a reversible leading-order swirl response for symmetric up-down forcing.

\paragraph{(iii) Swirl-Clock bookkeeping.} The SST Rosetta draft cleanly distinguishes the material scale
\begin{equation}
\label{eq:Omega0}
\Omega_0 \equiv \frac{\lVert \vswirl\rVert}{\rc},
\end{equation}
from the laboratory rotation rate $\Omega$, and retains the clock factor
\begin{equation}
\label{eq:SwirlClock}
\SwirlClock = \sqrt{1-\frac{v^2}{c^2}},
\qquad
\SwirlClockcw = \sqrt{1-\frac{v^2}{c^2}},
\end{equation}
as the canonical time-scaling variable. The two branch labels $\swirlarrow$ and $\swirlarrowcw$ are repurposed here as the internal branch index of the atomic sector.

\section{Minimal retained field sector}
The first retained sector is a transverse torsion/shear field. The minimal Lagrangian is
\begin{equation}
\label{eq:Ltorsion}
\mathcal{L}_{\rm torsion}
=
\frac12\rho_{\rm eff}\,\lVert \partial_t \mathbf{A}\rVert^2
-
\frac12 K\,\lVert \nabla\times\mathbf{A}\rVert^2,
\qquad
\nabla\cdot\mathbf{A}=0.
\end{equation}
Varying with respect to $\mathbf{A}$ gives
\begin{equation}
\label{eq:Awave-pre}
\rho_{\rm eff}\left(\partial_t^2\mathbf{A}+c^2\nabla\times\nabla\times\mathbf{A}\right)=0,
\qquad
c^2\equiv \frac{K}{\rho_{\rm eff}},
\end{equation}
which reduces, under $\nabla\cdot\mathbf{A}=0$, to
\begin{equation}
\label{eq:Awave}
\boxed{
\left(\frac{1}{c^2}\partial_t^2-\nabla^2\right)\mathbf{A}=0.
}
\end{equation}
This is the propagating field retained from the older torsion/photon sector. No magneto-optical or parity-odd vacuum extensions are retained in the present canonical core.

\section{Minimal retained source sector}
The second retained sector is the rotating-medium source law. The atomic theory proposed below does not require that every torsion excitation originate in a macroscopic rotating cylinder, but the rotating-cylinder calculation provides the cleanest explicit production mechanism for an SST torsion amplitude.

Define a dimensionless torsion-driving amplitude by normalizing the vertical vorticity to the canonical SST material rate~\eqref{eq:Omega0},
\begin{equation}
\label{eq:ThetaDef}
\boxed{
\Theta(\mathbf{r},t) \equiv \frac{\omega_z(\mathbf{r},t)}{\Omega_0},
\qquad
\Omega_0\equiv \frac{\lVert\vswirl\rVert}{\rc}.
}
\end{equation}
Equivalently, via~\eqref{eq:wzxi},
\begin{equation}
\label{eq:ThetaXi}
\boxed{
\Theta(\mathbf{r},t)=\frac{2\Omega}{\Omega_0}\,\partial_z\xi(\mathbf{r},t).
}
\end{equation}
Thus the source chain is
\begin{equation}
\label{eq:source-chain}
\xi \longrightarrow \omega_z \longrightarrow u_\theta \longrightarrow \Theta,
\end{equation}
where $\Theta$ is the quantity that couples the rotating-source sector to the branch-resolved atomic sector.

For a localized helical packet, a practical ansatz is
\begin{equation}
\label{eq:ThetaPacket}
\Theta(r,\phi,z,t)
=
A_\Theta
\exp\!\left(-\frac{r^2}{2w_r^2}\right)
\exp\!\left(-\frac{(z-vt)^2}{2w_z^2}\right)
\cos\!\left(m\phi+kz-\omega t\right).
\end{equation}
This profile is retained because it is the minimal object that carries amplitude, axial propagation, helical twist, and finite spatial support.

\section{Atomic sector: branch-resolved minimal Hamiltonian}
The atomic bridge model is built by retaining the ordinary Coulomb core and many-electron screening, and placing the novel SST content entirely in a branch-sensitive clock potential. The working Hamiltonian is
\begin{equation}
\label{eq:master-H}
\boxed{
\left[
-\frac{\hbar^2}{2\mu}\nabla^2
+V_{\rm core}(r)
+V_{\rm screen}[\rho](\mathbf r)
+V_{\rm clock}^{(\sigma)}[S_t]
\right]\Psi_{i,\sigma}
=\varepsilon_{i,\sigma}\Psi_{i,\sigma},
}
\end{equation}
with
\begin{equation}
\label{eq:Vcore}
V_{\rm core}(r)=-\frac{Ze^2}{4\pi\varepsilon_0 r}
\end{equation}
for the conservative baseline.

A provisional ``freeze v0.1'' removes the free branch couplings in favor of canonical SST scales. First, define the branch-coupling strength by the small dimensionless material ratio
\begin{equation}
\label{eq:eta0}
\boxed{
\eta_0 \equiv \frac{\lVert\vswirl\rVert}{c} = \frac{\alpha}{2} \approx 3.648676\times 10^{-3}.
}
\end{equation}
Then take the unperturbed atomic baseline to be branch-symmetric,
\begin{equation}
\label{eq:St0freeze}
S_{t,0}(r)\equiv 1,
\end{equation}
and define the two branch sectors by
\begin{equation}
\label{eq:StBranches}
\boxed{
S_t^{(\sigma)}(\mathbf r,t)=1+\sigma\,\eta_0\,\Theta(\mathbf r,t),
\qquad
\sigma=\pm 1,
}
\end{equation}
where $\sigma=+1$ corresponds to $\SwirlClock$ and $\sigma=-1$ to $\SwirlClockcw$.

Next, freeze the clock-energy scale to the Rydberg binding scale,
\begin{equation}
\label{eq:ERfreeze}
\boxed{
E_R \equiv \frac{\hbar^2}{2\mu (a_0^{\rm SST})^2}
= \frac{\mu e^4}{2(4\pi\varepsilon_0)^2\hbar^2}.
}
\end{equation}
The gradient coefficient is then fixed by the same atomic length,
\begin{equation}
\label{eq:lambda1freeze}
\boxed{
E_R (a_0^{\rm SST})^2 = \frac{\hbar^2}{2\mu}.
}
\end{equation}
The clock potential is therefore frozen to
\begin{equation}
\label{eq:Vclock}
\boxed{
V_{\rm clock}^{(\sigma)}
=
E_R\Bigl[1-\bigl(S_t^{(\sigma)}\bigr)^2 + (a_0^{\rm SST})^2\,\abs{\nabla S_t^{(\sigma)}}^2\Bigr].
}
\end{equation}
This choice eliminates the free parameters $(\chi,\lambda_0,\lambda_1)$ in favor of the already retained canonical scales $(\lVert\vswirl\rVert,\rc,\alpha,a_0^{\rm SST},E_R)$.

At first order in the torsion amplitude,
\begin{equation}
\label{eq:Vclock-linear}
\boxed{
\delta V_{\rm clock}^{(\sigma)}
\approx
-2\sigma E_R\eta_0\Theta,
}
\end{equation}
while the gradient contribution enters only at quadratic order,
\begin{equation}
\label{eq:Vclock-grad2}
\delta V_{\rm clock,grad}^{(\sigma)}
= E_R (a_0^{\rm SST})^2\eta_0^2\,\abs{\nabla\Theta}^2 + O(\Theta^3).
\end{equation}
Thus a torsion packet induces opposite branch shifts with no new fit constants.

\section{Many-electron closure and screening}
To reach shell structure, a many-electron mean-field closure is required. The electron density is defined by
\begin{equation}
\label{eq:rho-def}
\rho(\mathbf r)=\sum_{i,\sigma} f_{i,\sigma}\,\abs{\Psi_{i,\sigma}(\mathbf r)}^2,
\end{equation}
with occupation numbers $f_{i,\sigma}\in\{0,1\}$ at the effective level. The minimal retained screening functional is the Hartree term,
\begin{equation}
\label{eq:Vscreen}
\boxed{
V_{\rm screen}[\rho](\mathbf r)
=
\frac{e^2}{4\pi\varepsilon_0}
\int \frac{\rho(\mathbf r')}{\abs{\mathbf r-\mathbf r'}}\,d^3r'.
}
\end{equation}
No SST-specific correction to screening is introduced in the canonical starter model. This is deliberate: a predictive theory must first recover the ordinary many-electron limit before attempting branch- or medium-induced corrections to screening.

At this stage Pauli antisymmetry is retained as an effective low-energy rule rather than claimed as derived. The shell-state counts then remain
\begin{equation}
\label{eq:subshell-count}
N_{n\ell}=2(2\ell+1),
\end{equation}
and
\begin{equation}
\label{eq:shell-count}
N_n=\sum_{\ell=0}^{n-1}2(2\ell+1)=2n^2.
\end{equation}
In the present bridge model, the factor $2\ell+1$ arises from the ordinary rotational sectors $m_\ell=-\ell,\ldots,+\ell$, while the factor $2$ is interpreted as the low-energy branch doublet associated with $\SwirlClock$ and $\SwirlClockcw$.

\section{Hydrogenic limit and SST length scale}
The model must recover the ordinary hydrogenic spectrum when
\begin{equation}
V_{\rm screen}\to 0,
\qquad
\Theta\to 0,
\qquad
S_t^{(+)}=S_t^{(-)}.
\end{equation}
Then~\eqref{eq:master-H} reduces to the standard one-electron Coulomb problem,
\begin{equation}
\label{eq:En}
E_n=-\frac{\mu Z^2 e^4}{2(4\pi\varepsilon_0)^2\hbar^2}\frac{1}{n^2}.
\end{equation}
The SST atomic scale is tied to the canonical constants by
\begin{equation}
\label{eq:a0-sst}
\boxed{
a_0^{\rm SST}
=
\frac{F_{\rm swirl}^{\max}\,\rc^2}{M_e\,\lVert\vswirl\rVert^2}
=
\frac{2\rc}{\alpha^2}.
}
\end{equation}
Using the canonical values
\begin{equation}
\lVert\vswirl\rVert = 1.09384563\times 10^6\ \mathrm{m\,s^{-1}},
\quad
\rc = 1.40897017\times 10^{-15}\ \mathrm{m},
\quad
F_{\rm swirl}^{\max}=29.053507\ \mathrm{N},
\end{equation}
this gives
\begin{equation}
\label{eq:a0-num}
a_0^{\rm SST}
=
5.29177161\times 10^{-11}\ \mathrm{m},
\end{equation}
while the equivalent relation $2\rc/\alpha^2$ yields
\begin{equation}
\label{eq:a0-num2}
\frac{2\rc}{\alpha^2}=5.29177214\times 10^{-11}\ \mathrm{m},
\end{equation}
consistent with the ordinary Bohr scale to the expected rounding level.

This gives the hydrogenic SST radius law
\begin{equation}
\label{eq:RnZ}
R_{n,Z}^{\rm SST}\sim \frac{n^2}{Z}\,a_0^{\rm SST}.
\end{equation}
Thus the atomic bridge model retains the tested size law while reinterpreting $a_0$ as a medium-derived coherence scale. In the present draft this agreement is to be read as a \emph{consistency check on the retained canonical scales}, not as an independent prediction, since $(F_{\rm swirl}^{\max},\rc,\lVert\vswirl\rVert)$ were fixed upstream in the SST program.

\section{Transition amplitudes and torsion capture}
A torsion packet~\eqref{eq:ThetaPacket} drives atomic transitions through the clock-sector perturbation~\eqref{eq:Vclock-linear}. For an initial state $\ket{i,\sigma}$ and final state $\ket{f,\sigma'}$, the first-order matrix element is
\begin{equation}
\label{eq:Mfi}
M_{fi}^{\sigma'\sigma}=\matrixel{f,\sigma'}{\delta V_{\rm clock}}{i,\sigma}.
\end{equation}
The conservative resonance condition is
\begin{equation}
\label{eq:resonance}
\hbar\omega \approx \varepsilon_f-\varepsilon_i.
\end{equation}
If the packet carries azimuthal phase factor $e^{im\phi}$, then the orbital projection rule is expected to obey
\begin{equation}
\label{eq:dm}
\Delta m_\ell = m.
\end{equation}
Thus the torsion profile controls both the energy transfer and the transferred twist.

This suggests the following interpretation. Upward shell excitation under resonant incoming torsion is the atomic analogue of a ``giant jet,'' while downward relaxation accompanied by outgoing torsion radiation is the analogue of a ``sprite.'' These are analogical labels only; the retained mathematics is ordinary time-dependent perturbation theory acting on a branch-resolved Hamiltonian.

\section{Continuum sector and photoionization benchmark}
The sharpest next benchmark is photoionization rather than further shell bookkeeping. A branch-resolved atomic theory is incomplete until it specifies not only its bound spectrum but also its continuum and the dynamical channel by which an incoming torsion packet transfers energy into that continuum. The minimal conservative choice is to define the continuum as the positive-energy sector of the same branch-resolved Hamiltonian already introduced in Eq.~\eqref{eq:master-H}.

Write the static branch Hamiltonian as
\begin{equation}
\label{eq:H0sigma}
\boxed{
H_0^{(\sigma)}
=
-\frac{\hbar^2}{2\mu}\nabla^2
+
V_{\rm core}(r)
+
V_{\rm screen}[\rho](\mathbf r)
+
V_{{\rm clock},0}^{(\sigma)}(r),
\qquad
\sigma=\pm 1.
}
\end{equation}
The corresponding effective potential is
\begin{equation}
\label{eq:Veffsigma}
V_{\rm eff}^{(\sigma)}(\mathbf r)
=
V_{\rm core}(r)
+
V_{\rm screen}[\rho](\mathbf r)
+
V_{{\rm clock},0}^{(\sigma)}(r).
\end{equation}
Bound states satisfy
\begin{equation}
\label{eq:boundstates}
H_0^{(\sigma)}\psi_{n\ell m_\ell}^{(\sigma)}
=
E_{n\ell}^{(\sigma)}\psi_{n\ell m_\ell}^{(\sigma)},
\qquad
E_{n\ell}^{(\sigma)}<0,
\end{equation}
whereas continuum scattering states satisfy
\begin{equation}
\label{eq:continuumstates}
H_0^{(\sigma)}\psi_{\mathbf k}^{(\sigma)}
=
E_{\mathbf k}^{(\sigma)}\psi_{\mathbf k}^{(\sigma)},
\qquad
E_{\mathbf k}^{(\sigma)}>0.
\end{equation}
In the conservative asymptotic limit one imposes
\begin{equation}
\label{eq:clock-asymptotic}
V_{{\rm clock},0}^{(+)}(\infty)=V_{{\rm clock},0}^{(-)}(\infty)=0,
\end{equation}
so that the continuum threshold is branch-degenerate and may be taken as
\begin{equation}
\label{eq:continuum-threshold}
E_{\rm cont}^{(+)}(0)=E_{\rm cont}^{(-)}(0)=0.
\end{equation}
This is the simplest way to ensure that ionization is not put in by hand as a separate channel, but emerges as the $E>0$ part of the same atomic operator.

The next required step is to close the dynamical link between the propagating torsion field and the atomic driver $\Theta$. A minimal starter closure is to define $\Theta$ directly from the torsion/shear field through an axial projected curl,
\begin{equation}
\label{eq:ThetaA}
\boxed{
\Theta(\mathbf r,t)
=
\frac{1}{\Omega_0}
\hat{\mathbf z}\cdot(\nabla\times \mathbf A)(\mathbf r,t).
}
\end{equation}
This choice is not yet canonically derived from a deeper action, but it has three advantages: it makes $\Theta$ dimensionless, it ties the atomic driving amplitude to the retained torsion field rather than inserting it by hand, and it identifies the same helical packet profile with both a field configuration and an atomic perturbation source.

To permit back-reaction one must go one step further and source the field equation itself. The minimal branch-sensitive sourced wave equation is
\begin{equation}
\label{eq:Awave-source}
\boxed{
\left(\frac{1}{c^2}\partial_t^2-\nabla^2\right)\mathbf A
=
g_J \,\mathbf J_{\rm br},
}
\end{equation}
with branch current
\begin{equation}
\label{eq:Jbr}
\boxed{
\mathbf J_{\rm br}
=
\frac{\hbar}{2 i \mu}
\sum_{\sigma=\pm 1}
\sigma
\left(
\Psi_\sigma^*\nabla\Psi_\sigma
-
(\nabla\Psi_\sigma^*)\Psi_\sigma
\right).
}
\end{equation}
Equation~\eqref{eq:Jbr} should be read as a provisional closure ansatz rather than as a theorem of SST. Its purpose is to eliminate the prior ``stapled'' separation between the field and atomic sectors by making the field respond to branch-weighted atomic flow.

Once $\Theta$ is defined by Eq.~\eqref{eq:ThetaA}, the frozen clock coupling from Eq.~\eqref{eq:Vclock-linear} becomes a genuine interaction operator,
\begin{equation}
\label{eq:HintTheta}
\boxed{
\delta H_{\rm int}^{(\sigma)}(t)
=
-2\sigma E_R\eta_0\Theta(\mathbf r,t)
+O(\Theta^2).
}
\end{equation}
This is the weak-field quantum-absorption regime. In that regime the baseline photoelectric process is \emph{not} taken to be classical barrier suppression or ``branch shift larger than the binding energy.'' Instead, the incoming torsion packet transfers a discrete quantum $\hbar\omega$ into a bound state and promotes the electron into the continuum of the same Hamiltonian.

The branch-resolved ionization threshold from an initial bound state $i$ to a final continuum branch $\sigma_f$ is therefore
\begin{equation}
\label{eq:IonizationThresholdGeneral}
I_i^{(\sigma_i\to \sigma_f)}
=
E_{\rm cont}^{(\sigma_f)}(0)-E_i^{(\sigma_i)}.
\end{equation}
Under the asymptotic choice~\eqref{eq:continuum-threshold}, this reduces to
\begin{equation}
\label{eq:IonizationThresholdSimple}
I_i^{(\sigma_i)}=-E_i^{(\sigma_i)}.
\end{equation}
The photoionization condition is then frequency-thresholded,
\begin{equation}
\label{eq:PhotoThreshold}
\boxed{
\hbar\omega \ge I_i^{(\sigma_i\to \sigma_f)}.
}
\end{equation}
For an emitted free electron with asymptotic kinetic energy $K$, one obtains
\begin{equation}
\label{eq:KineticEnergyAtom}
\boxed{
K
=
\hbar\omega - I_i^{(\sigma_i\to \sigma_f)}.
}
\end{equation}
For a condensed system with work function $\phi^{(\sigma_i\to \sigma_f)}$, the analogous relation is
\begin{equation}
\label{eq:PhotoelectricSolid}
\boxed{
K_{\max}
=
\hbar\omega - \phi^{(\sigma_i\to \sigma_f)}.
}
\end{equation}
Thus the threshold is controlled by frequency, while the intensity enters only through the occupation number of the incoming torsion mode.

To show this explicitly, let $N_\omega$ denote the occupation number of an incoming torsion packet of frequency $\omega$. Then the continuum version of Fermi's golden rule becomes
\begin{equation}
\label{eq:GoldenRuleContinuum}
\boxed{
\Gamma_{i\to \mathbf k,\sigma_f}
=
\frac{2\pi}{\hbar}
N_\omega
\left|
\matrixel{\psi_{\mathbf k}^{(\sigma_f)}}{\delta H_{\rm int}^{(\sigma_f)}}{\psi_i^{(\sigma_i)}}
\right|^2
\delta\!\left(
E_{\mathbf k}^{(\sigma_f)}-E_i^{(\sigma_i)}-\hbar\omega
\right).
}
\end{equation}
This is the required rate structure if the SST torsion packet is to reproduce the key empirical logic of the photoelectric effect: an immediate threshold in frequency, and an emission rate proportional to flux or intensity rather than a threshold controlled by amplitude.

A useful distinction therefore emerges between two regimes. In the \emph{weak-field quantum absorption regime}, Eqs.~\eqref{eq:PhotoThreshold}--\eqref{eq:GoldenRuleContinuum} define the canonical baseline. In a possible later \emph{strong-field nonlinear regime}, sufficiently large coherent $\Theta$ could reshape the effective barrier and produce field-emission or over-the-barrier escape. That second regime should be treated as an extension and not conflated with the baseline photoelectric benchmark.

The branch-sensitive content enters through either the threshold or the matrix element. If the continuum is exactly branch-degenerate asymptotically, then branch dependence enters through $E_i^{(\sigma_i)}$ and through the matrix element in Eq.~\eqref{eq:GoldenRuleContinuum}. This permits a direct helicity/branch asymmetry observable,
\begin{equation}
\label{eq:Asymmetry}
\boxed{
\mathcal A(\omega)
=
\frac{\Gamma_{h=+1}-\Gamma_{h=-1}}{\Gamma_{h=+1}+\Gamma_{h=-1}},
}
\end{equation}
where $h=\pm 1$ labels the handedness of the incoming torsion packet. If $\mathcal A(\omega)=0$ identically once the couplings are frozen, then the present branch interpretation adds no atomic content beyond standard spin bookkeeping. If instead $\mathcal A(\omega)\neq 0$ with a definite symmetry and scale, then the theory acquires a genuine SST-specific test.

For this reason the photoionization benchmark is sharper than further shell counting. Shell capacities can be retained from effective low-energy quantum mechanics while Pauli is still deferred. By contrast, the continuum sector forces the theory to say how energy enters, what the final states are, and how the propagating field talks back to the atom. It is therefore the natural next benchmark for the retained canonical core.

\section{What is intentionally left outside the canonical core}
Three earlier ingredients are intentionally excluded from the present article.

First, the ultra-fast compressional branch with kinematic front speed $c_P\gg c$ is excluded. It is too loosely coupled to the retained atomic sector and too constrained by external bounds to serve as part of a first canonical core.

Second, the skyrmion-emission and topological-inheritance material is excluded from the canonical starter model. It is topologically interesting but not required to close the atomic Hamiltonian.

Third, vacuum magneto-optical extensions such as parity-odd Faraday rotation are excluded. They may become downstream consequences of the torsion field sector, but they are not needed to define the atomic bridge itself.

\section{Minimal predictive commitments}
The present model earns the label ``possibly canonical'' only if its commitments are kept narrow. The retained predictive commitments are:
\begin{enumerate}[label=(P\arabic*)]
    \item hydrogenic recovery: Eq.~\eqref{eq:En} must emerge when $\Theta\to 0$ and $V_{\rm screen}\to 0$;
    \item shell counting: the low-energy atomic sector must reproduce Eqs.~\eqref{eq:subshell-count} and~\eqref{eq:shell-count};
    \item branch symmetry: if $S_t^{(+)}=S_t^{(-)}$, the two branches must be degenerate;
    \item branch sensitivity: if $\Theta\neq 0$, then branch-sensitive line shifts or transition asymmetries are allowed and, in principle, measurable;
    \item frozen starter couplings: in the present v0.1 freeze, $\Theta$, $\eta_0$, and $V_{\rm clock}^{(\sigma)}$ are given explicitly by Eqs.~\eqref{eq:ThetaDef}, \eqref{eq:StBranches}, and~\eqref{eq:Vclock}, with no atom-by-atom retuning.
\end{enumerate}

The strongest falsifier is simple. If no branch-sensitive effects survive once the couplings are fixed by canon, then the branch interpretation reduces to a relabeling of ordinary spin without new atomic content. If branch-sensitive effects do survive, their magnitude and symmetry scaling become direct experimental targets.

\section{Status}
This article does not claim a first-principles derivation of fermionic exchange, nor a first-principles derivation of the branch-resolved clock sector beyond the provisional freeze given above. It is therefore not yet a closed atomic theory. What it does provide is a sharply reduced and internally consistent starter model:
\begin{equation}
\boxed{
\text{one propagating torsion field}
+\text{ one rotating-fluid source law}
+\text{ one branch-resolved atomic Hamiltonian}.
}
\end{equation}
That reduction is its main contribution. It turns three partially overlapping drafts into one narrow framework that can be tested, extended, or discarded term by term.

\paragraph{Analogy for a 10-year-old.} Imagine the atom as a globe with allowed storm patterns on it. A twisting wave packet can hit the globe and nudge one storm pattern into a higher allowed one. The two Swirl-Clock branches are like two opposite ways the storm can twist internally. The point of this paper is to say exactly which equations describe the wave, the nudge, and the storm pattern.

\begin{thebibliography}{99}

\bibitem{IskandaraniSST10}
O. Iskandarani,
\textit{Impulsive Axisymmetric Forcing in a Rotating Cylinder, Reversible Swirl Response, and Skyrmionic Photon Fields: Fluid Benchmarks and Fluid-Inspired Kinematic Hypotheses},
Zenodo preprint, 2026.
DOI: 10.5281/zenodo.17619170.

\bibitem{IskandaraniSST11}
O. Iskandarani,
\textit{Impulsive Axisymmetric Forcing in a Rotating Cylinder, Reversible Swirl Response, and Skyrmionic Photon Emission: Fluid Benchmarks and Fluid-Inspired Kinematic Hypotheses},
Zenodo preprint, 2026.
DOI: 10.5281/zenodo.17619573.

\bibitem{IskandaraniSST17}
O. Iskandarani,
\textit{Electromagnetism as Propagating Torsion in a Hydrodynamic Vacuum: A Geometric Unification via Cartan Structure Equations},
Zenodo preprint, 2026.
DOI: 10.5281/zenodo.17677074.

\bibitem{Greenspan1968}
H. P. Greenspan,
\textit{The Theory of Rotating Fluids},
Cambridge University Press, 1968.

\bibitem{Batchelor1967}
G. K. Batchelor,
\textit{An Introduction to Fluid Dynamics},
Cambridge University Press, 1967.

\bibitem{Vallis2017}
G. K. Vallis,
\textit{Atmospheric and Oceanic Fluid Dynamics},
2nd ed., Cambridge University Press, 2017.
DOI: 10.1017/9781107588417.

\bibitem{GriffithsSchroeter2018}
D. J. Griffiths and D. F. Schroeter,
\textit{Introduction to Quantum Mechanics},
3rd ed., Cambridge University Press, 2018.

\bibitem{BetheSalpeter1957}
H. A. Bethe and E. E. Salpeter,
\textit{Quantum Mechanics of One- and Two-Electron Atoms},
Springer, 1957.

\bibitem{Dirac1928}
P. A. M. Dirac,
``The Quantum Theory of the Electron,''
\textit{Proceedings of the Royal Society A} \textbf{117} (1928), 610--624.
DOI: 10.1098/rspa.1928.0023.

\bibitem{HehlObukhov2007}
F. W. Hehl and Y. N. Obukhov,
\textit{Elie Cartan's Torsion in Geometry and in Field Theory, an Essay},
Annales de la Fondation Louis de Broglie \textbf{32} (2007), 157--194.
Permalink: hal.science/hal-00408133.

\end{thebibliography}

\end{document}
